\documentclass[10pt,a4paper]{amsart}
\usepackage[margin=19mm]{geometry}
\usepackage{amsmath,amssymb,mathtools}
\usepackage[scr=rsfs,cal=euler]{mathalfa}
\usepackage{xfrac}
\usepackage[unicode,hidelinks,bookmarks=false]{hyperref}
\newcommand{\ie}{i.e.\ }
\newcommand{\cf}{cf.\ }
\newcommand{\resp}{respectively\ }
\newcommand{\Subsection}{\subsection}
\providecommand{\nobreakdash}{\nobreak\hbox{-}}
\setlength{\emergencystretch}{2em}
\multlinegap0pt
\usepackage{xcolor}
\usepackage{tikz}
\usepackage{enumerate}
\usepackage{amssymb}
\usepackage{mathtools}
\newtheorem{proposition}{Proposition}
\newcommand{\Q}{\mathbb{Q}}
\newcommand{\ord}{\operatorname{ord}}
\newcommand{\sg}{{}}
\newcommand{\sm}{{}}
\title{Computing $p$-adic heights on hyperelliptic curves}
\author{Stevan Gajovi\'c, J. Steffen M\"uller}
\date{Source: arXiv:2307.15787v4 (CC BY 4.0)}
\begin{document}
\maketitle

\noindent\emph{Source and licence.} Computing $p$-adic heights on hyperelliptic curves. Version arXiv:2307.15787v4 (CC BY 4.0), distributed by arXiv under the Creative Commons Attribution 4.0 International licence, http://creativecommons.org/licenses/by/4.0/, as recorded in the arXiv licence metadata for this exact version and shown on the abstract page https://arxiv.org/abs/2307.15787v4. Full licence notice: \texttt{licenses/arxiv-2307.15787-CC-BY-4.0.txt}.\par\smallskip

\noindent\textit{Editorial scope.} Counted unit: the article's proposition P:alphasecond (source0.tex lines 1195--1204). The article's complete original proof of it is delivered with it (source0.tex lines 1206--1250, one \texttt{proof} environment of 45 lines). No mathematics is written, repaired or re-derived: the statement and the proof are the article's own lines, in the article's own notation, and no retained line is substituted or re-flowed.  No further source-local context is needed: every cross-reference of every retained line resolves inside the two delivered spans.  Nothing was omitted from any retained span, so the package carries no omission record.  No retained line contains a citation, so the wrapper carries no bibliography and no bibliography entry.  No retained line contains a figure environment, an image-inclusion command or drawing code, so no image is delivered and the package declares no asset dependency.  Editorial formatting only, in addition: an AMS \texttt{amsart} preamble that restates the article's own declarations and macros used by the retained lines, namely the environments \texttt{proposition} on separate counters, together with the standard packages those lines need.  The retained environments print in this document as Proposition 1 in order of appearance, whereas the article prints the same items with the numbering of its own full document.  The complete native source of the article as distributed in the arXiv e-print for this exact version is bundled unchanged as \texttt{sources/144916/source0.tex}.
\begin{proposition}\label{P:alphasecond}\hfill
\begin{itemize}
\item[(a)] The differential $\alpha$ is holomorphic at both $\infty_{\pm}$.
\item[(b)] \sg{The
  differential $\alpha$ is essentially of the second kind. More precisely,
    $\alpha$ is holomorphic at non-Weierstrass points, and has  residue~0
    along the discs $D_P$ for all Weierstrass points $P$.} 
    %along a sufficiently small annulus containing Weierstrass points is zero.}
\end{itemize} 
\end{proposition}

\begin{proof}
To prove (a) we only consider one point at infinity, say $\infty_-$; the
  proof is the same for the other one. The action of Frobenius on the
  uniformizer $t=\frac{1}{x}$ is 
\[
\phi^* t=\phi^*\left(\dfrac{1}{x}\right)=\dfrac{1}{x^p}=t^p\,.
\]
We know that $\dfrac{x^gdx}{y}=\dfrac{dt}{t}+A(t) dt $, where $A(t)\in \Q_p[[t]]$.
Hence
\[
\phi^*\left(\dfrac{x^gdx}{y}\right)=p\dfrac{dt}{t}+pt^{t-1}A(t^p)dt\,.
\]
It follows that
\[
\alpha=(pt^{p-1}A(t^p)-pA(t))dt
\]
is holomorphic at $\infty_-$. 

For (b), we first show that $\alpha$ is holomorphic at non-Weierstrass
  points. We already know that \sg{$\omega_g$} is holomorphic at every point
  $P=(a,b)\in C(\overline{\Q}_p)$ with $b\neq 0$. The function $t=x-a$ is a uniformizer at $P$, thus $dx=dt$. We compute 
\begin{equation}\label{eq:Frob-of-wg}
 \sg{\phi^*(\omega_g)}=\phi^*\left(\dfrac{x^gdx}{y}\right)=\dfrac{px^{pg+p-1}dx}{y^p}\displaystyle\sum_{i\geq
  0}\dbinom{-\frac{1}{2}}{k}\dfrac{(f(x^p)-f(x)^p)^i}{y^{2pi}}\,,    
\end{equation}
so that  \sg{$\ord_P\frac{\phi^*(\omega_g)}{dt}=\ord_P\frac{\phi^*(\omega_g)}{dx}\geq 0$}.
Hence the $t$-adic valuation of $\alpha$ is nonnegative.

A uniformizer at a Weierstrass point $P=(a,0)$ is given by $t=y$. Then $x$ is an
  even function of $t$ as we now explain. Let $x=\sum_{n\geq 0} e_nt^n$ for
  some $e_n\in\overline{\Q}_p$, then, inductively equating powers of $t$ on
  both sides of the equation $t^2=f(x)$, we see that $f(e_0)=0$, and then for odd $n$, $f'(e_0)e_n=0$. Since $f'(e_0)\neq 0$, it follows that $e_n=0$ for odd $n$.
We now prove that the same holds for the expansion of $\alpha$ in $t$. Using the relation
$$\dfrac{dx}{y}=\dfrac{2dy}{f'(x)}$$
it is clear that 
$$p\dfrac{x^gdx}{y}=p\dfrac{2x^gdy}{f'(x)}$$
has only even powers of $t$. For the other term of $\alpha$, we
  re-express~\eqref{eq:Frob-of-wg} using $dy$:
$$\phi^*\left(\dfrac{x^gdx}{y}\right)=\dfrac{2px^{pg+p-1}dy}{f'(x)y^{p-1}}\displaystyle\sum_{i\geq 0}\dbinom{-\frac{1}{2}}{k}\dfrac{(f(x^p)-f(x)^p)^i}{y^{2pi}},$$
from which we see that all powers of $y$ are even. Hence there are no odd
  powers of $t$ in the expansion of $\alpha$ and therefore \sm{$\alpha$ has
  trivial residue along $D_P$. }
  %\sg{if this is correct:
  %the residue along any disc containing $P$ is zero, so} (b) follows.
\end{proof}

\end{document}
